\section{Experimental Setup}
\label{sec:experimental-setup}

\subsection{Hardware and Software}

Experiments were conducted on the following platform:

\begin{table}[h]
\centering
\caption{Benchmarking environment}
\label{tab:environment}
\begin{tabular}{ll}
\toprule
\textbf{Component} & \textbf{Specification} \\
\midrule
OS & Windows 11 Pro (amd64) \\
CPU & 20 logical cores \\
Go version & 1.25.7 \\
gnark-crypto & v0.20.1 \\
BN254 field & $p = 2^{254} + \epsilon$ \\
\bottomrule
\end{tabular}
\end{table}

The gnark-crypto library~\cite{gnark2024crypto} provides
assembly-optimized BN254 field arithmetic for AMD64 using
Montgomery multiplication. The \texttt{fr.Element} type stores
field elements in 4$\times$64-bit limb representation, avoiding
heap allocations that plague \texttt{math/big.Int} implementations.

\subsection{Implementation Details}

\subsubsection{Poseidon2 Hash Function}

The Poseidon2 permutation is configured with width 2 (2-to-1
compression), $r_f = 6$ full rounds, and $r_p = 50$ partial rounds
over the BN254 scalar field. The compression function computes:

\begin{equation}
H(a, b) = \text{Perm}(a, b)[1] + b
\label{eq:poseidon2}
\end{equation}

where $\text{Perm}$ is the Poseidon2 permutation and element
index 1 of the output is used as the digest, following the
feed-forward construction. A global permutation instance is
shared across all tree operations to amortize initialization cost.

\subsubsection{Sparse Merkle Tree}

The SMT implementation uses a map-based in-memory store keyed by
$(level, path)$ pairs. Key features:

\begin{itemize}
  \item \textbf{Precomputed defaults:} An array of $d + 1$ default
    hashes is computed at initialization, where
    $\text{default}[i] = H(\text{default}[i-1], \text{default}[i-1])$.
  \item \textbf{Leaf construction:}
    $\text{leaf} = H(\text{key}, \text{value})$.
  \item \textbf{Path traversal:} Bit $i$ of the key determines
    left (0) or right (1) child at level $i$.
  \item \textbf{Sparse optimization:} Non-existent siblings use
    precomputed default hashes, avoiding storage of empty subtrees.
\end{itemize}

\subsubsection{Optimized Variant}

The optimized implementation replaces all \texttt{big.Int} operations
with gnark-crypto's \texttt{fr.Element}:

\begin{lstlisting}[language=Go, caption=Optimized SMT node storage]
type SMTOptimized struct {
  depth    int
  nodes    map[string]fr.Element
  defaults []fr.Element
  perm     hash.Hash
}
\end{lstlisting}

Key differences from the baseline:
\begin{itemize}
  \item Field elements stored in Montgomery form (no conversion
    overhead).
  \item Pre-allocated permutation buffer reused across hash calls.
  \item \texttt{InsertUint64} fast path for keys representable in
    64 bits (depth $\leq 64$).
  \item Proof operations (\texttt{GetProofUint64},
    \texttt{VerifyProofUint64}) operate entirely in
    \texttt{fr.Element} domain.
\end{itemize}

\subsection{Benchmark Protocol}

Each benchmark follows this protocol:

\begin{enumerate}
  \item Initialize tree with target parameters (depth, entry count).
  \item Execute 10 warmup iterations (results discarded).
  \item Execute 50 measured iterations with distinct random seeds.
  \item Record wall-clock time via \texttt{time.Now()} /
    \texttt{time.Since()}.
  \item Compute statistical summary: mean, standard deviation,
    min, max, P50, P95, P99.
  \item Force garbage collection via \texttt{runtime.GC()} before
    memory measurements.
  \item Export results to JSON for analysis.
\end{enumerate}

Key generation uses deterministic sequential integers to ensure
reproducibility across runs and comparability with the TypeScript
RU-V1 baseline, which used the same key generation strategy.
